\section{Core Concepts}
\label{thesis:intro-concepts}

\begin{figure}[t]
    \centering
    \includegraphics[width=0.6\columnwidth]{sw-stack.drawio.pdf}
    \caption[Hardware and software stack in a modern computing system]{Overview of the hardware and software stack in a modern computing system. The security of each layer depends on the layers below, meaning that a compromise at a certain layer might break the confidentiality and integrity not only in that layer, but also all the layers above.}
    \label{fig:phd-intro-stack}
\end{figure}

In this section, we provide a short summary of basic principles of computing and
security that are relevant for this dissertation. Here, we do not aim to be
comprehensive, and we make a few simplifications for the sake of brevity.

\cref{fig:phd-intro-stack} shows the main hardware and software layers and
components of a modern computing system. At the bottom, the hardware layer
includes the main CPU and firmware, memory and disks, network interfaces, and
peripheral devices such as clock, co-processors, GPU, and so on. Here, we may
also find security modules such as the \ac{TPM}~\cite{tpmdoc}. Going up there is
the virtualization layer, where a hypervisor or \ac{VMM} is responsible for
managing one or more \acp{VM}. Physical resources such as CPU cores and memory
are then \emph{virtualized} and can be assigned to multiple \acp{VM}, enabling
concurrency. Note that the virtualization layer is optional, though it is widely
used in cloud settings to run multiple workloads on the same physical node,
effectively optimizing resource utilization and reducing costs. Next there is
the OS layer, which includes components such as the kernel, kernel modules and
device drivers. This layer is responsible for bridging the gap between
applications and the underlying hardware resources (which may be virtualized if
running within a \ac{VM}), and for providing isolation, communication and
scheduling mechanisms. Lastly, the application layer contains applications,
libraries and runtimes that run the actual workloads.

In a traditional system, the security of a certain layer is enforced by the
layer below, providing mechanisms such as resource separation and virtualization
to isolate different applications (enforced by the OS layer) and \acp{VM}
(enforced by the virtualization layer). This also means that lower layers in the
stack have full control on the upper layers and can access all their code and
data. Therefore, a compromise at a certain layer might break the confidentiality
and integrity both in that layer and all the ones above. To this end, the term
\ac{TCB} is defined as the sum of all hardware and software components that are
critical to guarantee the confidentiality and integrity of a certain
application. To build robust systems, it is therefore important to keep the
\ac{TCB} as small as possible.